\chapter{Nhớ và Quên}
\markboth{Nhớ và Quên}{Remembering and Forgetting}

\section{Nhớ ơn để sống đàng hoàng hơn.}
\begin{truyen}{Nhớ ơn để sống đàng hoàng hơn.}{Gratitude keeps character awake.}
\chuhoa{M}{ỗi lần đi ngang cổng trường cũ, anh lại nhớ cô giáo từng ở lại thêm giờ để kèm mình. Sự nhớ ấy không làm anh yếu mềm; nó giữ cho anh không cư xử hỗn với người đi sau. Có những ký ức không để ta luyến tiếc, mà để ta tử tế hơn.}
\textit{(Each time he passed his old school gate, he remembered the teacher who had stayed late to tutor him. That memory did not make him sentimental; it kept him from becoming rude to those who came after him. Some memories are not for nostalgia, but for decency.)}
\end{truyen}
\begin{baihoc}
Nhớ đúng điều đáng nhớ giúp con người không cạn nghĩa.
\textit{(Remembering what deserves remembrance keeps a person from becoming empty of gratitude.)}
\end{baihoc}
\begin{ghinhoanh}
\textbf{Short English to remember:}

Gratitude is memory with character.
\end{ghinhoanh}
\begin{tuhocnhanh}
1. Ai là người em cần nhớ ơn rõ hơn? \textit{(Whom do you need to remember with greater gratitude?)}

2. Ký ức biết ơn nào đang giúp em sống tốt hơn? \textit{(Which grateful memory is helping you live better?)}
\end{tuhocnhanh}
\ngancach

\section{Quên một nỗi đau để không sống mãi trong cay đắng.}
\begin{truyen}{Quên một nỗi đau để không sống mãi trong cay đắng.}{Some forgetting is healing.}
\chuhoa{C}{ô không quên bài học từ lần bị phản bội, nhưng cô tập quên dần cảm giác muốn trả đũa. Nếu giữ mãi cơn đau, nó sẽ xin ở luôn trong mình. Quên ở đây không phải là ngu ngơ, mà là không cho điều xấu quyền cai trị phần đời còn lại.}
\textit{(She did not forget the lesson of betrayal, but she tried to let go of the urge for revenge. If pain is held forever, it asks to live inside us permanently. Forgetting here is not foolishness; it is refusing to let the wrong rule the rest of life.)}
\end{truyen}
\begin{baihoc}
Có những điều cần nhớ để khôn, nhưng cũng có những phần cần quên để lành.
\textit{(Some things must be remembered to become wise, yet other parts must be forgotten to become whole.)}
\end{baihoc}
\begin{ghinhoanh}
\textbf{Short English to remember:}

Remember the lesson, not the poison.
\end{ghinhoanh}
\begin{tuhocnhanh}
1. Điều gì em nên giữ như bài học nhưng không nên giữ như vết thương? \textit{(What should you keep as a lesson but not as an open wound?)}

2. Em có đang sống quá lâu trong một nỗi đau cũ không? \textit{(Are you living too long inside an old pain?)}
\end{tuhocnhanh}
\ngancach

\section{Nhớ lỗi người khác mà quên lỗi mình.}
\begin{truyen}{Nhớ lỗi người khác mà quên lỗi mình.}{Selective memory feeds self-righteousness.}
\chuhoa{A}{nh kể rất rành những lần người khác làm mình tổn thương, nhưng lại rất mơ hồ khi nhắc tới các lỗi mình từng gây ra. Cách nhớ lệch ấy khiến anh thấy mình luôn là nạn nhân và khó trưởng thành. Nhớ không công bằng là một dạng tự đánh lừa bản thân.}
\textit{(He remembered clearly every time others hurt him, yet became strangely vague about the harm he had caused. Such uneven memory allowed him to see himself only as a victim and made growth difficult. Unfair remembering is a subtle form of self-deception.)}
\end{truyen}
\begin{baihoc}
Muốn trưởng thành, phải nhớ cả phần mình cần sửa.
\textit{(To grow, we must remember the parts of ourselves that need correction.)}
\end{baihoc}
\begin{ghinhoanh}
\textbf{Short English to remember:}

Fair memory includes yourself.
\end{ghinhoanh}
\begin{tuhocnhanh}
1. Em có đang nhớ thiên lệch theo hướng có lợi cho mình không? \textit{(Are you remembering things in a way that always favors yourself?)}

2. Lỗi nào của em cần được nhìn lại công bằng hơn? \textit{(Which of your mistakes needs to be faced more honestly?)}
\end{tuhocnhanh}
\ngancach

\section{Quên lời hứa nhỏ làm mòn lòng tin lớn.}
\begin{truyen}{Quên lời hứa nhỏ làm mòn lòng tin lớn.}{Small forgotten promises weaken big trust.}
\chuhoa{M}{ột người cha hứa đón con rồi quên vì bận. Chuyện ấy không quá lớn với người lớn, nhưng với đứa trẻ, đó là cảm giác bị bỏ quên. Có những điều nhỏ trong trí nhớ của ta lại rất lớn trong trái tim của người khác.}
\textit{(A father promised to pick up his child and forgot because of work. To an adult it seemed small, but to the child it felt like abandonment. Some things that seem minor in our memory are huge in another person's heart.)}
\end{truyen}
\begin{baihoc}
Quên việc nhỏ nhiều lần có thể làm hỏng một sự tin cậy lớn.
\textit{(Repeatedly forgetting small things can damage a larger trust.)}
\end{baihoc}
\begin{ghinhoanh}
\textbf{Short English to remember:}

Small promises carry big trust.
\end{ghinhoanh}
\begin{tuhocnhanh}
1. Em có đang coi nhẹ lời hứa với ai không? \textit{(Are you treating a promise to someone too lightly?)}

2. Điều nhỏ nào em cần nhớ kỹ hơn để giữ lòng tin? \textit{(What small thing should you remember more carefully to protect trust?)}
\end{tuhocnhanh}
\ngancach

\section{Nhớ mình từng yếu để bớt khắc nghiệt với người khác.}
\begin{truyen}{Nhớ mình từng yếu để bớt khắc nghiệt với người khác.}{Memory can soften judgment.}
\chuhoa{K}{hi thấy học sinh chậm tiến, cô giáo nhớ lại chính mình từng sợ Toán, từng ngại phát biểu, từng đứng trước lớp mà run. Sự nhớ ấy làm cô không dễ mắng nặng. Ai còn nhớ những ngày mình yếu thường ít dùng sức mạnh để chèn ép người khác.}
\textit{(When she saw slow students, the teacher remembered her own fear of math, her hesitation to speak, and her trembling before class. That memory kept her from harsh scolding. Those who remember their weaker days are less likely to misuse strength.)}
\end{truyen}
\begin{baihoc}
Ký ức về sự yếu đuối có thể trở thành nguồn của lòng cảm thông.
\textit{(Memory of weakness can become a source of compassion.)}
\end{baihoc}
\begin{ghinhoanh}
\textbf{Short English to remember:}

Remembering weakness can grow compassion.
\end{ghinhoanh}
\begin{tuhocnhanh}
1. Quãng khó nào của em có thể giúp em hiểu người khác hơn? \textit{(Which hard season of yours can help you understand others better?)}

2. Em cần bớt khắt khe với ai bằng cách nhớ lại điều gì? \textit{(With whom do you need to be less harsh by remembering something from your own past?)}
\end{tuhocnhanh}
\ngancach
